% Evaluation section of the paper. Meant to be included from paper.tex with
% % Evaluation section of the paper. Meant to be included from paper.tex with
% % Evaluation section of the paper. Meant to be included from paper.tex with
% % Evaluation section of the paper. Meant to be included from paper.tex with
% \input{evaluation} at the point in the main text where the evaluation belongs.
\section{Evaluating the annotation pipeline}
\label{sec:evaluation}

\subsection{Gold standard and models compared}

We evaluate the pipeline against a single hand-annotated novel,
\emph{L'Ange et les pervers} (Lucie Delarue-Mardrus, 1930), for which a human annotator
followed our annotation codebook (the same one given to the models) to produce 14 gold
relationships. This is a genuine limitation of the present evaluation: one novel gives us
a first read of how the pipeline behaves, not a corpus-level estimate of its accuracy, and
the numbers below should be read as a pilot, not as a final benchmark. Extending the gold
standard to further novels of the corpus is future work.

Five large language models, all served through OpenRouter with the same system prompt,
codebook, and JSON schema, were run on this novel (Table~\ref{tab:models}).

\begin{table}[h]
  \centering
  \small
  \begin{tabular}{ll}
    \toprule
    Provider & Model \\
    \midrule
    DeepSeek  & V4 Pro \\
    Anthropic & Claude Sonnet 5 \\
    Anthropic & Claude Opus 5 \\
    OpenAI    & GPT-5.6 Luna Pro \\
    OpenAI    & GPT-5.6 Terra \\
    \bottomrule
  \end{tabular}
  \caption{The five models compared, all called through OpenRouter with an identical
  system prompt, codebook, and JSON schema.}
  \label{tab:models}
\end{table}

\subsection{Two capacities, two metrics}

Annotating a novel asks two different things of a model: deciding \emph{which}
relationships exist, and \emph{describing} each one along the eleven closed fields of the
codebook (gender configuration, structure, status, and so on). A model can find exactly
the right relationships and mislabel all of them, or the reverse, so we score these two
capacities separately before combining them into a single number.

For a novel, let $G = \{g_1, \dots, g_n\}$ be the set of gold relationships and
$P = \{p_1, \dots, p_m\}$ the set of relationships predicted by a model. Every relationship
$r$ carries a participant set $C(r)$ and a label vector $\ell(r) = (\ell_1(r), \dots, \ell_K(r))$
over the codebook's closed fields $\mathcal{F}$, grouped in three families:

\begin{itemize}
  \item \emph{participants and shape}: \texttt{relation\_type}, \texttt{structure};
  \item \emph{nature and rules of the bond}: \texttt{status}, \texttt{reciprocity},
    \texttt{exclusivity}, \texttt{adultery};
  \item \emph{representation and narrative role}: \texttt{explicitness},
    \texttt{explicitness\_in\_plot}, \texttt{centrality}, \texttt{framing}, \texttt{outcome}.
\end{itemize}

This gives $K = 11$ fields in $\mathcal{F}$. The gold file used below predates
\texttt{explicitness\_in\_plot}, added in schema v2, so this pilot scores grid agreement over
the remaining $K = 10$ fields; a twelfth field, \texttt{confidence}, is logged by the model but
is metadata rather than part of $\mathcal{F}$, and is not scored. $G$ and $P$ are unordered
sets with no shared index, so identification requires an explicit alignment step before any
field can be compared.

\paragraph{Identification.} Each relationship is keyed by its (canonicalised) participant
set, $\kappa(r) = C(r)$ — e.g.\ $\kappa(r) = \texttt{laurette\_aimee}$ for a relationship
between Laurette and Aimée. The matching set

\begin{align}
M = \{(g,p) \in G \times P \;:\; \kappa(g) = \kappa(p)\}
\end{align}

is computed as an exact, one-to-one hash join on this key: no similarity function and no
threshold enter the identification score, which is a deterministic function of the two
annotation files. From $M$ we compute precision, recall, and their harmonic mean:

\begin{align}
P_{\mathrm{det}} &= \frac{|M|}{|P|} , \\
R_{\mathrm{det}} &= \frac{|M|}{|G|} , \\
F_{1,\mathrm{det}} &= \frac{2\,P_{\mathrm{det}}\,R_{\mathrm{det}}}{P_{\mathrm{det}} + R_{\mathrm{det}}} .
\end{align}

$F_{1,\mathrm{det}}$ only asks whether the model found the right relationships; it says
nothing about whether it described them correctly.

\paragraph{Grid agreement.} For every matched pair $(g,p) \in M$, we compare the closed
fields of $\mathcal{F}$ one by one — \texttt{relation\_type}, \texttt{structure},
\texttt{status}, \texttt{reciprocity}, \texttt{exclusivity}, \texttt{adultery},
\texttt{explicitness}, \texttt{centrality}, \texttt{framing}, \texttt{outcome} for this pilot,
per the field count noted above — and average the agreement with a uniform weight $w_k = 1/K$,
so every field counts the same regardless of which of the three families it belongs to:

\begin{align}
s_k(\ell_k(g), \ell_k(p)) &= \mathbb{1}[\ell_k(g) = \ell_k(p)] , \\
a(g,p) &= \frac{1}{K}\sum_{k \in \mathcal{F}} s_k(\ell_k(g), \ell_k(p)) , \\
\bar{A} &= \frac{1}{|M|}\sum_{(g,p)\in M} a(g,p) .
\end{align}

Here $s_k$ is the strict (exact-match) credit function. $\bar A$ is only computed on matched
relationships: it measures description quality independently of how many relationships were
found.

\paragraph{Composite score.} Precision and recall re-weighted by how well each matched
pair was described,

\begin{align}
P_{\mathrm{soft}} &= \frac{\sum_{(g,p) \in M} a(g,p)}{|P|} , \\
R_{\mathrm{soft}} &= \frac{\sum_{(g,p) \in M} a(g,p)}{|G|} ,
\end{align}

combine into a single composite score,

\begin{align}
\mathrm{sF}_1 = \frac{2\,P_{\mathrm{soft}}\,R_{\mathrm{soft}}}{P_{\mathrm{soft}} + R_{\mathrm{soft}}} = \bar{A} \cdot F_{1,\mathrm{det}} ,
\end{align}

the second equality following directly from substituting $P_{\mathrm{soft}} = \bar A \, P_{\mathrm{det}}$
and $R_{\mathrm{soft}} = \bar A \, R_{\mathrm{det}}$ into the harmonic mean, so that $\bar A$
cancels once from the denominator and survives once in the numerator. $\mathrm{sF}_1$
therefore decomposes exactly into how many of the right relationships were found times how
well the found ones were described, with $\mathrm{sF}_1 \le F_{1,\mathrm{det}}$ and equality
only if every matched pair is described perfectly. We report the triple
$(F_{1,\mathrm{det}}, \bar A, \mathrm{sF}_1)$ rather than $\mathrm{sF}_1$ alone, since two
models can reach the same composite score for very different reasons.

\subsection{Results}

\begin{table}[h]
  \centering
  \small
  \begin{tabular}{lccccc}
    \toprule
    Model & $P_{\mathrm{det}}$ & $R_{\mathrm{det}}$ & $F_{1,\mathrm{det}}$ & $\bar{A}$ & $\mathrm{sF}_1$ \\
    \midrule
    \textbf{Claude Opus 5} & 0.71 & 0.71 & 0.71 & 0.67 & \textbf{0.48} \\
    GPT-5.6 Terra          & 1.00 & 0.43 & 0.60 & 0.73 & 0.44 \\
    Claude Sonnet 5        & 1.00 & 0.43 & 0.60 & 0.67 & 0.40 \\
    GPT-5.6 Luna Pro       & 1.00 & 0.36 & 0.53 & 0.60 & 0.32 \\
    DeepSeek V4 Pro        & 0.86 & 0.43 & 0.57 & 0.47 & 0.27 \\
    \bottomrule
  \end{tabular}
  \caption{Identification, grid agreement, and composite score for the five models
  benchmarked on the single annotated novel, sorted by $\mathrm{sF}_1$.}
  \label{tab:model-results}
\end{table}

Claude Opus 5 recovers exactly seven tenths of the gold relationships, and recovers them
evenly (precision and recall are both 0.71); together with a solid grid agreement
($\bar A = 0.67$) this gives it the best composite score. The other four models are all
perfectly precise (every relationship they propose is a real one, $P_{\mathrm{det}} = 1$),
but each finds notably fewer of the fourteen gold relationships, which caps their recall
and their composite scores below Claude Opus 5. Description quality among the models that
do find relationships is fairly close ($\bar A$ between 0.47 and 0.73); none of the five
gets every field right on any matched relationship ($\bar A = 1$ never occurs at the level
of a full relationship in this pilot).

On the strength of this pilot, \textbf{Claude Opus 5} is the model we carry forward for
annotating the rest of the corpus, while keeping in mind that this ranking rests on a
single novel and 14 gold relationships, and should be revisited once further novels are
added to the gold standard.

 at the point in the main text where the evaluation belongs.
\section{Evaluating the annotation pipeline}
\label{sec:evaluation}

\subsection{Gold standard and models compared}

We evaluate the pipeline against a single hand-annotated novel,
\emph{L'Ange et les pervers} (Lucie Delarue-Mardrus, 1930), for which a human annotator
followed our annotation codebook (the same one given to the models) to produce 14 gold
relationships. This is a genuine limitation of the present evaluation: one novel gives us
a first read of how the pipeline behaves, not a corpus-level estimate of its accuracy, and
the numbers below should be read as a pilot, not as a final benchmark. Extending the gold
standard to further novels of the corpus is future work.

Five large language models, all served through OpenRouter with the same system prompt,
codebook, and JSON schema, were run on this novel (Table~\ref{tab:models}).

\begin{table}[h]
  \centering
  \small
  \begin{tabular}{ll}
    \toprule
    Provider & Model \\
    \midrule
    DeepSeek  & V4 Pro \\
    Anthropic & Claude Sonnet 5 \\
    Anthropic & Claude Opus 5 \\
    OpenAI    & GPT-5.6 Luna Pro \\
    OpenAI    & GPT-5.6 Terra \\
    \bottomrule
  \end{tabular}
  \caption{The five models compared, all called through OpenRouter with an identical
  system prompt, codebook, and JSON schema.}
  \label{tab:models}
\end{table}

\subsection{Two capacities, two metrics}

Annotating a novel asks two different things of a model: deciding \emph{which}
relationships exist, and \emph{describing} each one along the eleven closed fields of the
codebook (gender configuration, structure, status, and so on). A model can find exactly
the right relationships and mislabel all of them, or the reverse, so we score these two
capacities separately before combining them into a single number.

For a novel, let $G = \{g_1, \dots, g_n\}$ be the set of gold relationships and
$P = \{p_1, \dots, p_m\}$ the set of relationships predicted by a model. Every relationship
$r$ carries a participant set $C(r)$ and a label vector $\ell(r) = (\ell_1(r), \dots, \ell_K(r))$
over the codebook's closed fields $\mathcal{F}$, grouped in three families:

\begin{itemize}
  \item \emph{participants and shape}: \texttt{relation\_type}, \texttt{structure};
  \item \emph{nature and rules of the bond}: \texttt{status}, \texttt{reciprocity},
    \texttt{exclusivity}, \texttt{adultery};
  \item \emph{representation and narrative role}: \texttt{explicitness},
    \texttt{explicitness\_in\_plot}, \texttt{centrality}, \texttt{framing}, \texttt{outcome}.
\end{itemize}

This gives $K = 11$ fields in $\mathcal{F}$. The gold file used below predates
\texttt{explicitness\_in\_plot}, added in schema v2, so this pilot scores grid agreement over
the remaining $K = 10$ fields; a twelfth field, \texttt{confidence}, is logged by the model but
is metadata rather than part of $\mathcal{F}$, and is not scored. $G$ and $P$ are unordered
sets with no shared index, so identification requires an explicit alignment step before any
field can be compared.

\paragraph{Identification.} Each relationship is keyed by its (canonicalised) participant
set, $\kappa(r) = C(r)$ — e.g.\ $\kappa(r) = \texttt{laurette\_aimee}$ for a relationship
between Laurette and Aimée. The matching set

\begin{align}
M = \{(g,p) \in G \times P \;:\; \kappa(g) = \kappa(p)\}
\end{align}

is computed as an exact, one-to-one hash join on this key: no similarity function and no
threshold enter the identification score, which is a deterministic function of the two
annotation files. From $M$ we compute precision, recall, and their harmonic mean:

\begin{align}
P_{\mathrm{det}} &= \frac{|M|}{|P|} , \\
R_{\mathrm{det}} &= \frac{|M|}{|G|} , \\
F_{1,\mathrm{det}} &= \frac{2\,P_{\mathrm{det}}\,R_{\mathrm{det}}}{P_{\mathrm{det}} + R_{\mathrm{det}}} .
\end{align}

$F_{1,\mathrm{det}}$ only asks whether the model found the right relationships; it says
nothing about whether it described them correctly.

\paragraph{Grid agreement.} For every matched pair $(g,p) \in M$, we compare the closed
fields of $\mathcal{F}$ one by one — \texttt{relation\_type}, \texttt{structure},
\texttt{status}, \texttt{reciprocity}, \texttt{exclusivity}, \texttt{adultery},
\texttt{explicitness}, \texttt{centrality}, \texttt{framing}, \texttt{outcome} for this pilot,
per the field count noted above — and average the agreement with a uniform weight $w_k = 1/K$,
so every field counts the same regardless of which of the three families it belongs to:

\begin{align}
s_k(\ell_k(g), \ell_k(p)) &= \mathbb{1}[\ell_k(g) = \ell_k(p)] , \\
a(g,p) &= \frac{1}{K}\sum_{k \in \mathcal{F}} s_k(\ell_k(g), \ell_k(p)) , \\
\bar{A} &= \frac{1}{|M|}\sum_{(g,p)\in M} a(g,p) .
\end{align}

Here $s_k$ is the strict (exact-match) credit function. $\bar A$ is only computed on matched
relationships: it measures description quality independently of how many relationships were
found.

\paragraph{Composite score.} Precision and recall re-weighted by how well each matched
pair was described,

\begin{align}
P_{\mathrm{soft}} &= \frac{\sum_{(g,p) \in M} a(g,p)}{|P|} , \\
R_{\mathrm{soft}} &= \frac{\sum_{(g,p) \in M} a(g,p)}{|G|} ,
\end{align}

combine into a single composite score,

\begin{align}
\mathrm{sF}_1 = \frac{2\,P_{\mathrm{soft}}\,R_{\mathrm{soft}}}{P_{\mathrm{soft}} + R_{\mathrm{soft}}} = \bar{A} \cdot F_{1,\mathrm{det}} ,
\end{align}

the second equality following directly from substituting $P_{\mathrm{soft}} = \bar A \, P_{\mathrm{det}}$
and $R_{\mathrm{soft}} = \bar A \, R_{\mathrm{det}}$ into the harmonic mean, so that $\bar A$
cancels once from the denominator and survives once in the numerator. $\mathrm{sF}_1$
therefore decomposes exactly into how many of the right relationships were found times how
well the found ones were described, with $\mathrm{sF}_1 \le F_{1,\mathrm{det}}$ and equality
only if every matched pair is described perfectly. We report the triple
$(F_{1,\mathrm{det}}, \bar A, \mathrm{sF}_1)$ rather than $\mathrm{sF}_1$ alone, since two
models can reach the same composite score for very different reasons.

\subsection{Results}

\begin{table}[h]
  \centering
  \small
  \begin{tabular}{lccccc}
    \toprule
    Model & $P_{\mathrm{det}}$ & $R_{\mathrm{det}}$ & $F_{1,\mathrm{det}}$ & $\bar{A}$ & $\mathrm{sF}_1$ \\
    \midrule
    \textbf{Claude Opus 5} & 0.71 & 0.71 & 0.71 & 0.67 & \textbf{0.48} \\
    GPT-5.6 Terra          & 1.00 & 0.43 & 0.60 & 0.73 & 0.44 \\
    Claude Sonnet 5        & 1.00 & 0.43 & 0.60 & 0.67 & 0.40 \\
    GPT-5.6 Luna Pro       & 1.00 & 0.36 & 0.53 & 0.60 & 0.32 \\
    DeepSeek V4 Pro        & 0.86 & 0.43 & 0.57 & 0.47 & 0.27 \\
    \bottomrule
  \end{tabular}
  \caption{Identification, grid agreement, and composite score for the five models
  benchmarked on the single annotated novel, sorted by $\mathrm{sF}_1$.}
  \label{tab:model-results}
\end{table}

Claude Opus 5 recovers exactly seven tenths of the gold relationships, and recovers them
evenly (precision and recall are both 0.71); together with a solid grid agreement
($\bar A = 0.67$) this gives it the best composite score. The other four models are all
perfectly precise (every relationship they propose is a real one, $P_{\mathrm{det}} = 1$),
but each finds notably fewer of the fourteen gold relationships, which caps their recall
and their composite scores below Claude Opus 5. Description quality among the models that
do find relationships is fairly close ($\bar A$ between 0.47 and 0.73); none of the five
gets every field right on any matched relationship ($\bar A = 1$ never occurs at the level
of a full relationship in this pilot).

On the strength of this pilot, \textbf{Claude Opus 5} is the model we carry forward for
annotating the rest of the corpus, while keeping in mind that this ranking rests on a
single novel and 14 gold relationships, and should be revisited once further novels are
added to the gold standard.

 at the point in the main text where the evaluation belongs.
\section{Evaluating the annotation pipeline}
\label{sec:evaluation}

\subsection{Gold standard and models compared}

We evaluate the pipeline against a single hand-annotated novel,
\emph{L'Ange et les pervers} (Lucie Delarue-Mardrus, 1930), for which a human annotator
followed our annotation codebook (the same one given to the models) to produce 14 gold
relationships. This is a genuine limitation of the present evaluation: one novel gives us
a first read of how the pipeline behaves, not a corpus-level estimate of its accuracy, and
the numbers below should be read as a pilot, not as a final benchmark. Extending the gold
standard to further novels of the corpus is future work.

Five large language models, all served through OpenRouter with the same system prompt,
codebook, and JSON schema, were run on this novel (Table~\ref{tab:models}).

\begin{table}[h]
  \centering
  \small
  \begin{tabular}{ll}
    \toprule
    Provider & Model \\
    \midrule
    DeepSeek  & V4 Pro \\
    Anthropic & Claude Sonnet 5 \\
    Anthropic & Claude Opus 5 \\
    OpenAI    & GPT-5.6 Luna Pro \\
    OpenAI    & GPT-5.6 Terra \\
    \bottomrule
  \end{tabular}
  \caption{The five models compared, all called through OpenRouter with an identical
  system prompt, codebook, and JSON schema.}
  \label{tab:models}
\end{table}

\subsection{Two capacities, two metrics}

Annotating a novel asks two different things of a model: deciding \emph{which}
relationships exist, and \emph{describing} each one along the eleven closed fields of the
codebook (gender configuration, structure, status, and so on). A model can find exactly
the right relationships and mislabel all of them, or the reverse, so we score these two
capacities separately before combining them into a single number.

For a novel, let $G = \{g_1, \dots, g_n\}$ be the set of gold relationships and
$P = \{p_1, \dots, p_m\}$ the set of relationships predicted by a model. Every relationship
$r$ carries a participant set $C(r)$ and a label vector $\ell(r) = (\ell_1(r), \dots, \ell_K(r))$
over the codebook's closed fields $\mathcal{F}$, grouped in three families:

\begin{itemize}
  \item \emph{participants and shape}: \texttt{relation\_type}, \texttt{structure};
  \item \emph{nature and rules of the bond}: \texttt{status}, \texttt{reciprocity},
    \texttt{exclusivity}, \texttt{adultery};
  \item \emph{representation and narrative role}: \texttt{explicitness},
    \texttt{explicitness\_in\_plot}, \texttt{centrality}, \texttt{framing}, \texttt{outcome}.
\end{itemize}

This gives $K = 11$ fields in $\mathcal{F}$. The gold file used below predates
\texttt{explicitness\_in\_plot}, added in schema v2, so this pilot scores grid agreement over
the remaining $K = 10$ fields; a twelfth field, \texttt{confidence}, is logged by the model but
is metadata rather than part of $\mathcal{F}$, and is not scored. $G$ and $P$ are unordered
sets with no shared index, so identification requires an explicit alignment step before any
field can be compared.

\paragraph{Identification.} Each relationship is keyed by its (canonicalised) participant
set, $\kappa(r) = C(r)$ — e.g.\ $\kappa(r) = \texttt{laurette\_aimee}$ for a relationship
between Laurette and Aimée. The matching set

\begin{align}
M = \{(g,p) \in G \times P \;:\; \kappa(g) = \kappa(p)\}
\end{align}

is computed as an exact, one-to-one hash join on this key: no similarity function and no
threshold enter the identification score, which is a deterministic function of the two
annotation files. From $M$ we compute precision, recall, and their harmonic mean:

\begin{align}
P_{\mathrm{det}} &= \frac{|M|}{|P|} , \\
R_{\mathrm{det}} &= \frac{|M|}{|G|} , \\
F_{1,\mathrm{det}} &= \frac{2\,P_{\mathrm{det}}\,R_{\mathrm{det}}}{P_{\mathrm{det}} + R_{\mathrm{det}}} .
\end{align}

$F_{1,\mathrm{det}}$ only asks whether the model found the right relationships; it says
nothing about whether it described them correctly.

\paragraph{Grid agreement.} For every matched pair $(g,p) \in M$, we compare the closed
fields of $\mathcal{F}$ one by one — \texttt{relation\_type}, \texttt{structure},
\texttt{status}, \texttt{reciprocity}, \texttt{exclusivity}, \texttt{adultery},
\texttt{explicitness}, \texttt{centrality}, \texttt{framing}, \texttt{outcome} for this pilot,
per the field count noted above — and average the agreement with a uniform weight $w_k = 1/K$,
so every field counts the same regardless of which of the three families it belongs to:

\begin{align}
s_k(\ell_k(g), \ell_k(p)) &= \mathbb{1}[\ell_k(g) = \ell_k(p)] , \\
a(g,p) &= \frac{1}{K}\sum_{k \in \mathcal{F}} s_k(\ell_k(g), \ell_k(p)) , \\
\bar{A} &= \frac{1}{|M|}\sum_{(g,p)\in M} a(g,p) .
\end{align}

Here $s_k$ is the strict (exact-match) credit function. $\bar A$ is only computed on matched
relationships: it measures description quality independently of how many relationships were
found.

\paragraph{Composite score.} Precision and recall re-weighted by how well each matched
pair was described,

\begin{align}
P_{\mathrm{soft}} &= \frac{\sum_{(g,p) \in M} a(g,p)}{|P|} , \\
R_{\mathrm{soft}} &= \frac{\sum_{(g,p) \in M} a(g,p)}{|G|} ,
\end{align}

combine into a single composite score,

\begin{align}
\mathrm{sF}_1 = \frac{2\,P_{\mathrm{soft}}\,R_{\mathrm{soft}}}{P_{\mathrm{soft}} + R_{\mathrm{soft}}} = \bar{A} \cdot F_{1,\mathrm{det}} ,
\end{align}

the second equality following directly from substituting $P_{\mathrm{soft}} = \bar A \, P_{\mathrm{det}}$
and $R_{\mathrm{soft}} = \bar A \, R_{\mathrm{det}}$ into the harmonic mean, so that $\bar A$
cancels once from the denominator and survives once in the numerator. $\mathrm{sF}_1$
therefore decomposes exactly into how many of the right relationships were found times how
well the found ones were described, with $\mathrm{sF}_1 \le F_{1,\mathrm{det}}$ and equality
only if every matched pair is described perfectly. We report the triple
$(F_{1,\mathrm{det}}, \bar A, \mathrm{sF}_1)$ rather than $\mathrm{sF}_1$ alone, since two
models can reach the same composite score for very different reasons.

\subsection{Results}

\begin{table}[h]
  \centering
  \small
  \begin{tabular}{lccccc}
    \toprule
    Model & $P_{\mathrm{det}}$ & $R_{\mathrm{det}}$ & $F_{1,\mathrm{det}}$ & $\bar{A}$ & $\mathrm{sF}_1$ \\
    \midrule
    \textbf{Claude Opus 5} & 0.71 & 0.71 & 0.71 & 0.67 & \textbf{0.48} \\
    GPT-5.6 Terra          & 1.00 & 0.43 & 0.60 & 0.73 & 0.44 \\
    Claude Sonnet 5        & 1.00 & 0.43 & 0.60 & 0.67 & 0.40 \\
    GPT-5.6 Luna Pro       & 1.00 & 0.36 & 0.53 & 0.60 & 0.32 \\
    DeepSeek V4 Pro        & 0.86 & 0.43 & 0.57 & 0.47 & 0.27 \\
    \bottomrule
  \end{tabular}
  \caption{Identification, grid agreement, and composite score for the five models
  benchmarked on the single annotated novel, sorted by $\mathrm{sF}_1$.}
  \label{tab:model-results}
\end{table}

Claude Opus 5 recovers exactly seven tenths of the gold relationships, and recovers them
evenly (precision and recall are both 0.71); together with a solid grid agreement
($\bar A = 0.67$) this gives it the best composite score. The other four models are all
perfectly precise (every relationship they propose is a real one, $P_{\mathrm{det}} = 1$),
but each finds notably fewer of the fourteen gold relationships, which caps their recall
and their composite scores below Claude Opus 5. Description quality among the models that
do find relationships is fairly close ($\bar A$ between 0.47 and 0.73); none of the five
gets every field right on any matched relationship ($\bar A = 1$ never occurs at the level
of a full relationship in this pilot).

On the strength of this pilot, \textbf{Claude Opus 5} is the model we carry forward for
annotating the rest of the corpus, while keeping in mind that this ranking rests on a
single novel and 14 gold relationships, and should be revisited once further novels are
added to the gold standard.

 at the point in the main text where the evaluation belongs.
\section{Evaluating the annotation pipeline}
\label{sec:evaluation}

\subsection{Gold standard and models compared}

We evaluate the pipeline against a single hand-annotated novel,
\emph{L'Ange et les pervers} (Lucie Delarue-Mardrus, 1930), for which a human annotator
followed our annotation codebook (the same one given to the models) to produce 14 gold
relationships. This is a genuine limitation of the present evaluation: one novel gives us
a first read of how the pipeline behaves, not a corpus-level estimate of its accuracy, and
the numbers below should be read as a pilot, not as a final benchmark. Extending the gold
standard to further novels of the corpus is future work.

Five large language models, all served through OpenRouter with the same system prompt,
codebook, and JSON schema, were run on this novel (Table~\ref{tab:models}).

\begin{table}[h]
  \centering
  \small
  \begin{tabular}{ll}
    \toprule
    Provider & Model \\
    \midrule
    DeepSeek  & V4 Pro \\
    Anthropic & Claude Sonnet 5 \\
    Anthropic & Claude Opus 5 \\
    OpenAI    & GPT-5.6 Luna Pro \\
    OpenAI    & GPT-5.6 Terra \\
    \bottomrule
  \end{tabular}
  \caption{The five models compared, all called through OpenRouter with an identical
  system prompt, codebook, and JSON schema.}
  \label{tab:models}
\end{table}

\subsection{Two capacities, two metrics}

Annotating a novel asks two different things of a model: deciding \emph{which}
relationships exist, and \emph{describing} each one along the eleven closed fields of the
codebook (gender configuration, structure, status, and so on). A model can find exactly
the right relationships and mislabel all of them, or the reverse, so we score these two
capacities separately before combining them into a single number.

For a novel, let $G = \{g_1, \dots, g_n\}$ be the set of gold relationships and
$P = \{p_1, \dots, p_m\}$ the set of relationships predicted by a model. Every relationship
$r$ carries a participant set $C(r)$ and a label vector $\ell(r) = (\ell_1(r), \dots, \ell_K(r))$
over the codebook's closed fields $\mathcal{F}$, grouped in three families:

\begin{itemize}
  \item \emph{participants and shape}: \texttt{relation\_type}, \texttt{structure};
  \item \emph{nature and rules of the bond}: \texttt{status}, \texttt{reciprocity},
    \texttt{exclusivity}, \texttt{adultery};
  \item \emph{representation and narrative role}: \texttt{explicitness},
    \texttt{explicitness\_in\_plot}, \texttt{centrality}, \texttt{framing}, \texttt{outcome}.
\end{itemize}

This gives $K = 11$ fields in $\mathcal{F}$. The gold file used below predates
\texttt{explicitness\_in\_plot}, added in schema v2, so this pilot scores grid agreement over
the remaining $K = 10$ fields; a twelfth field, \texttt{confidence}, is logged by the model but
is metadata rather than part of $\mathcal{F}$, and is not scored. $G$ and $P$ are unordered
sets with no shared index, so identification requires an explicit alignment step before any
field can be compared.

\paragraph{Identification.} Each relationship is keyed by its (canonicalised) participant
set, $\kappa(r) = C(r)$ — e.g.\ $\kappa(r) = \texttt{laurette\_aimee}$ for a relationship
between Laurette and Aimée. The matching set

\begin{align}
M = \{(g,p) \in G \times P \;:\; \kappa(g) = \kappa(p)\}
\end{align}

is computed as an exact, one-to-one hash join on this key: no similarity function and no
threshold enter the identification score, which is a deterministic function of the two
annotation files. From $M$ we compute precision, recall, and their harmonic mean:

\begin{align}
P_{\mathrm{det}} &= \frac{|M|}{|P|} , \\
R_{\mathrm{det}} &= \frac{|M|}{|G|} , \\
F_{1,\mathrm{det}} &= \frac{2\,P_{\mathrm{det}}\,R_{\mathrm{det}}}{P_{\mathrm{det}} + R_{\mathrm{det}}} .
\end{align}

$F_{1,\mathrm{det}}$ only asks whether the model found the right relationships; it says
nothing about whether it described them correctly.

\paragraph{Grid agreement.} For every matched pair $(g,p) \in M$, we compare the closed
fields of $\mathcal{F}$ one by one — \texttt{relation\_type}, \texttt{structure},
\texttt{status}, \texttt{reciprocity}, \texttt{exclusivity}, \texttt{adultery},
\texttt{explicitness}, \texttt{centrality}, \texttt{framing}, \texttt{outcome} for this pilot,
per the field count noted above — and average the agreement with a uniform weight $w_k = 1/K$,
so every field counts the same regardless of which of the three families it belongs to:

\begin{align}
s_k(\ell_k(g), \ell_k(p)) &= \mathbb{1}[\ell_k(g) = \ell_k(p)] , \\
a(g,p) &= \frac{1}{K}\sum_{k \in \mathcal{F}} s_k(\ell_k(g), \ell_k(p)) , \\
\bar{A} &= \frac{1}{|M|}\sum_{(g,p)\in M} a(g,p) .
\end{align}

Here $s_k$ is the strict (exact-match) credit function. $\bar A$ is only computed on matched
relationships: it measures description quality independently of how many relationships were
found.

\paragraph{Composite score.} Precision and recall re-weighted by how well each matched
pair was described,

\begin{align}
P_{\mathrm{soft}} &= \frac{\sum_{(g,p) \in M} a(g,p)}{|P|} , \\
R_{\mathrm{soft}} &= \frac{\sum_{(g,p) \in M} a(g,p)}{|G|} ,
\end{align}

combine into a single composite score,

\begin{align}
\mathrm{sF}_1 = \frac{2\,P_{\mathrm{soft}}\,R_{\mathrm{soft}}}{P_{\mathrm{soft}} + R_{\mathrm{soft}}} = \bar{A} \cdot F_{1,\mathrm{det}} ,
\end{align}

the second equality following directly from substituting $P_{\mathrm{soft}} = \bar A \, P_{\mathrm{det}}$
and $R_{\mathrm{soft}} = \bar A \, R_{\mathrm{det}}$ into the harmonic mean, so that $\bar A$
cancels once from the denominator and survives once in the numerator. $\mathrm{sF}_1$
therefore decomposes exactly into how many of the right relationships were found times how
well the found ones were described, with $\mathrm{sF}_1 \le F_{1,\mathrm{det}}$ and equality
only if every matched pair is described perfectly. We report the triple
$(F_{1,\mathrm{det}}, \bar A, \mathrm{sF}_1)$ rather than $\mathrm{sF}_1$ alone, since two
models can reach the same composite score for very different reasons.

\subsection{Results}

\begin{table}[h]
  \centering
  \small
  \begin{tabular}{lccccc}
    \toprule
    Model & $P_{\mathrm{det}}$ & $R_{\mathrm{det}}$ & $F_{1,\mathrm{det}}$ & $\bar{A}$ & $\mathrm{sF}_1$ \\
    \midrule
    \textbf{Claude Opus 5} & 0.71 & 0.71 & 0.71 & 0.67 & \textbf{0.48} \\
    GPT-5.6 Terra          & 1.00 & 0.43 & 0.60 & 0.73 & 0.44 \\
    Claude Sonnet 5        & 1.00 & 0.43 & 0.60 & 0.67 & 0.40 \\
    GPT-5.6 Luna Pro       & 1.00 & 0.36 & 0.53 & 0.60 & 0.32 \\
    DeepSeek V4 Pro        & 0.86 & 0.43 & 0.57 & 0.47 & 0.27 \\
    \bottomrule
  \end{tabular}
  \caption{Identification, grid agreement, and composite score for the five models
  benchmarked on the single annotated novel, sorted by $\mathrm{sF}_1$.}
  \label{tab:model-results}
\end{table}

Claude Opus 5 recovers exactly seven tenths of the gold relationships, and recovers them
evenly (precision and recall are both 0.71); together with a solid grid agreement
($\bar A = 0.67$) this gives it the best composite score. The other four models are all
perfectly precise (every relationship they propose is a real one, $P_{\mathrm{det}} = 1$),
but each finds notably fewer of the fourteen gold relationships, which caps their recall
and their composite scores below Claude Opus 5. Description quality among the models that
do find relationships is fairly close ($\bar A$ between 0.47 and 0.73); none of the five
gets every field right on any matched relationship ($\bar A = 1$ never occurs at the level
of a full relationship in this pilot).

On the strength of this pilot, \textbf{Claude Opus 5} is the model we carry forward for
annotating the rest of the corpus, while keeping in mind that this ranking rests on a
single novel and 14 gold relationships, and should be revisited once further novels are
added to the gold standard.

